\documentclass[11pt,letter,oneside]{article}
\usepackage{sectsty,setspace} 
\usepackage[top=1.00in, bottom=1.0in, left=1in, right=1in]{geometry} 
\usepackage{graphicx}
\usepackage{epstopdf}
\usepackage{amsmath,latexsym,amssymb,wasysym}
\usepackage[english]{babel}
\usepackage{cite}
\usepackage{citesupernumber}
\usepackage{citecollapse}
\usepackage{wrapfig}
\usepackage{hyperref}
\usepackage{float}
\restylefloat{table}
\usepackage[table]{xcolor} % http://ctan.org/pkg/xcolor
\usepackage{gensymb} % You have to have this to use \degree

\setlength\parindent{0pt} % no indents throughout

\setlength{\abovecaptionskip}{1pt}
\setlength{\belowcaptionskip}{-10pt}
\usepackage[font={small}]{caption}

\parskip=6pt
\pagenumbering{arabic}
\pagestyle{plain}
% squeeze spacey
\linespread{0.99}
\addtolength{\itemsep}{-0.05in}
\usepackage{multicol}
 
\newenvironment{smitemize}{
\begin{itemize}
  \setlength{\itemsep}{1pt}
  \setlength{\parskip}{0pt}
  \setlength{\parsep}{0pt}}
{\end{itemize}
}

\usepackage{fancyhdr}
\pagestyle{fancy}
\fancyhead[LO]{Career}
\fancyhead[RO]{E. M. Wolkovich}

\newcommand{\Section}[1]{\vspace{-8pt}\section{\hskip -1em.~~#1}\vspace{-3pt}} 


\begin{document}
\bibliographystyle{/Users/Lizzie/Documents/EndnoteRelated/Bibtex/styles/ajhg}
\renewcommand{\refname}{\CHead{}}
\begin{centering}
{\bf Career trajectory}\\
\end{centering}
\vspace{1.5ex}
\doublespacing
I started my career fascinated by how humans have altered space. After millennia of habitat change by people, the massive modifications of the 20th century had led to the rise of spatial ecology---a discipline that tasked itself with understanding and predicting the consequences of habitat loss and fragmentation for plant and animal communities. I joined the charge. For my PhD I studied plant invasions in San Diego, where a system of canyons too steep to build on, and parks provided an ideal test system of `islands' and `mainlands' of habitat embedded in a sprawling urban matrix.

Yet, by the end of my PhD, I was less certain that it was the space that mattered, and more convinced that it was time. I found the biggest impacts of invasive grasses came through their timing. They responded more rapidly to rainfall, growing sooner and faster than native species, then senescing earlier, causing a cascade of ecosystem changes that promoted their growth and spread. 

After my PhD I thus focused on phenology---the timing of recurring life history events---and the most reported biological impact of climate change. I found that what I noticed during my PhD appeared across other systems: non-native and invasive species from the Northeast to Central Plains of North America had distinct phenologies from native species. More and more I began to notice how much the changes in communities and ecosystems were shaped by time---through altered disturbance cycles, introduced species, and climate change, but I found little theory to guide me. Fundamentally, we lack a way to predict the consequences when human actions alter biological time. 

Focused on this problem, from my postdoctoral period to today, I have worked to build the field of temporal ecology, which studies how time structures populations, communities and ecosystems. For several years I turned from field work to ecological bioinformatics, leveraging complex datasets from herbaria, historical texts and other long-term records. I led an interdisciplinary working group on plant phenology (competitively funded by NCEAS, 2010-2012 and I was funded by an NSF Postdoctoral Research Fellowship in Biology from 2010-2011). We compiled phenological data on \(>\)5,000 plant species with \(>\)30 years of time-series data on approximately 2,000 of those species, creating the most species-rich database available, and now the major citation for the Intergovernmental Panel on Climate Change (and many others) for current estimates of plant phenological shifts with climate change (and cited over 1000 times). With the data we showed that species exhibit a wide but predictable set of responses to temperature and precipitation across habitats, from the tropics, to the temperate zone where warmer winter temperatures may delay flowering and where sensitivity to temperature is greatest for early-flowering species. 

After this, I started the Temporal Ecology Lab at Harvard University, moving to the University of British Columbia (UBC) in 2018 where I continued work on my predictive framework. My lab's approach leverages theory from community ecology with interdisciplinary perspectives from evolutionary biology, climatology, engineering and data science to identify the mechanisms underlying shifts in species and communities. We combine field and lab experiments, gradient studies and synthesis of short and long-term data with bespoke sensors, AI and Bayesian models. Together we aim to elevate global change ecology to forecast beyond trends---to predict variation and complex feedbacks---by building from fundamental temperature response curves to community assembly and forecasting of shifts in communities and ecosystems. Over this time, I have earned many national and international awards, including an Early Career Fellowship from the Ecological Society of America, and a Canada Research Chair in Temporal Ecology.

My lab has focused on how biological time has been reshaped by anthropogenic warming. We have showed that major variation due to the complexity of measuring `time' as climate change accelerates biological responses. We found that dozens of studies finding that sensitivity to temperature has declined with warming can be predicted from the non-linearity of temperature responses. Relatedly, we also showed that experiments under-predicted  plant responses to temperature, due mainly to unrealistic shifts in soil moisture in experiments. Since then, we have developed a unified temperature response model to estimate variation across species and sub-species from observational data to predict both variability and future trends. 

Understanding how both environmental and genetic variation drive thermal responses is critical to accurate forecasts with continued climate change, especially for forest and agricultural systems. My lab was the first to estimate the importance of genetic diversity in heat tolerance to mitigating crop losses with climate change---showing that shifting to more heat-tolerant, late-ripening varieties of winegrapes could halve projected losses, saving a third of the world's growing regions. I am routinely asked to speak about this work as it has become a pivotal citation in discussions of adapting systems to continued warming.

My lab now aims to build on these broad results to understand how longer seasons impact tree growth. Theoretically more time to grow should lead to more growth, but it does not, at least not consistently. My work to date has outlined the major drivers of this variability (see \emph{Statement of Plans}), but progress will require an interdisciplinary approach. To that aim, my lab is working to leverage new modeling methods from statistics, advances in sensor technology from engineering, improved downscaled forecasts from climatology, and integrating evolutionary history to better understand species-level variation in response to warming. My team recently developed a new image capture machine to rapidly measure tree rings (paper in review by Fong \emph{et al.}), and we have active work underway sampling trees across elevational gradients at Mount Rainier National Park, where the growing season shortens considerably as you move higher. We have also started work on better ways to forecast. Models we recently developed found an almost equal role of the environment and evolutionary history in predicting future responses, and my proposed work builds on this by integrating species-specific responses to help understand when and why trees grow more with longer seasons. 


\end{document}


% At UBC my lab has worked together on global meta-analyses to better understand the current state of knowledge and gaps and in two major systems to test hypotheses and models: winegrapes and forest trees. We leverage the depth of physiological knowledge and extensive data resources of winegrapes to develop cutting-edge models that provide answers to critical questions in how best to adapt agriculture to climate change. In particular, we focus on different winegrape varieties which are distinct genotypes, such as Pinot noir or Merlot, that vary tremendously in their phenology, and resilience to drought and heat extremes. Much of this knowledge we then use in our study of forest tree communities to understand how species, communities, and carbon storage across North American forests may shift with warming. 



% Through my work to date, I have become a leader in two parallel tracks of research needed for my predictive framework: developing and testing new theory of community assembly via phenology (reviewed in 54), and advancing understanding of the fundamental cues that shape phenological events (36, 48) and the sequence of events each year (37, 51). While my work on the former has focused especially on forest tree species, I have worked on both trees and winegrapes for the latter. I have found my work on winegrapes has been especially valuable both because the depth of data and knowledge advances my understanding of natural communities, and because of the direct and clear impacts. Recently, we developed phenological models of winegrapes from budburst to maturity to show that if growers shift varieties of winegrapes with warming, it could save up to half of all winegrowing regions (43); this research provides a major example of how the intersection of biological diversity and human adaptation will define the resilience of many ecosystems with climate change. 

% My proposed research here extends my previous work, but with a sharp focus on building better temperature response models for winegrapes---with implications for diverse plant responses to climate change. Phenology is effectively a metric of biological time, thus my project is focused on our need for better models of time. When I started my research career 20 years ago, I was interested in space---its influence had become obvious because humans had so altered it. With climate change, we see how massively humans have warped time, and we are challenged to better understand and predict its consequences. 